\section{The wedge frame and the stabilized corner}\label{sec:corner}

\subsection{The complete quadratic frame}
For real vectors $x,y$ indexed by an ordered set of labels, write
$z_{ab}=x_ay_b-x_by_a$ for labels $a<b$.

\begin{lemma}[The complete quadratic frame]\label{lem:frame}
Let $F(x,y)$ be a real homogeneous quartic that vanishes whenever
$x=0$, $y=0$, or $y=x$. If $F$ is a sum of squares of real homogeneous
quadratics, then every square is a linear combination of the $z_{ab}$.
Moreover the kernel of the map $Q\mapsto z^TQz$ on symmetric matrices
is exactly the span of the Pl\"ucker matrices
\[
 z^TP_{abcd}z=2(z_{ab}z_{cd}-z_{ac}z_{bd}+z_{ad}z_{bc}),
 \qquad a<b<c<d.
\]
In particular, if $F=z^TGz$, then
\begin{equation}\label{eq:complete-frame}
 F\text{ is SOS}\quad\Longleftrightarrow\quad
 G+E\succeq0\text{ for some }E\in\Span\{P_{abcd}\}.
\end{equation}
\end{lemma}
\begin{proof}
Write $F=\sum_jq_j^2$. At each point with $y=0$, the sum of real
squares is zero; hence every $q_j(x,0)$ is identically zero. The same
holds at $x=0$. A homogeneous quadratic thus contains only terms
$x_ay_b$, so $q_j=x^TA_jy$. Now $q_j(x,x)=0$ for all real $x$ implies
$A_j+A_j^T=0$. Therefore
$q_j=\sum_{a<b}(A_j)_{ab}z_{ab}$. The coefficient vectors of these
quadratics give a positive semidefinite Gram matrix.

For the relation space, group products of wedges according to their
multiset of labels. Polynomial monomials from different multisets
cannot cancel. A multiset with a repeated label has at most one
nonzero pairing, up to the antisymmetric signs; its wedge product is a
nonzero polynomial. On four distinct labels there are precisely the
three products
$z_{ab}z_{cd},z_{ac}z_{bd},z_{ad}z_{bc}$.
Their displayed alternating combination is zero. Any two are linearly
independent: for example their coefficients of
$x_ax_cy_by_d$ and $x_ax_by_cy_d$ already have rank two.
Thus their relation space is one-dimensional, generated by the
Pl\"ucker identity. This proves the complete kernel description.

If an SOS Gram represents the same polynomial as $G$, their
difference belongs to this kernel. Conversely a positive semidefinite
representative has a real Gram factorization, whose rows are the
required wedge quadratics. This proves \eqref{eq:complete-frame}.
\end{proof}

The form $F_N$ and each of its restrictions vanish when $x=0$, $y=0$ or
$x=y$, so Lemma~\ref{lem:frame} applies to them. Positivity in
\eqref{eq:complete-frame} is required on all of wedge space, not only on
the decomposable vectors $x\wedge y$; this is the gap described in
\S\ref{sec:mech-wedges}.

\subsection{Coordinates and conventions}
Let $N\ge2m$. For $0\le p<m$ the retained symbol pairs at depth $p$ are
\[
 w_p^-=w_{\lambda_p^-},\qquad w_p^+=w_{\lambda_p^+},\qquad
 \lambda_p^\pm=\pm(N-1-p),
\]
where $w_\lambda=(x_\lambda,y_\lambda)$; the diagonal at depth $p$ from
either end has $p+1$ entries. Their wedges are the pure wedges
\(u_{pq}=w_p^-\wedge w_q^-\) and \(v_{pq}=w_p^+\wedge w_q^+\) for \(p<q\),
and the mixed wedges \(z^{\rm mix}_{pq}=w_p^-\wedge w_q^+\) for all \(p,q\).
The notation of the analytic argument is collected here.
\begin{center}
\begin{tabular}{@{}p{.24\textwidth}p{.72\textwidth}@{}}
\toprule
Symbol & Meaning\\\midrule
$w_\lambda$; $w_p^\pm$ & symbol pair at label $\lambda$; retained pair at depth $p$\\
$u_{pq},v_{pq}$; $z^{\rm mix}_{pq}$ & pure wedges of the two sides; mixed wedges\\
$D,K,B$ & fixed pure diagonal, pure cross and mixed matrices \eqref{eq:corner-baselines}\\
$E$, $\Theta$; $\Psi=M-B$ & pure four-form, pure cross correction; mixed correction\\
$P_\pm$, $M$ & averaged pure and mixed blocks \eqref{eq:corner-sectors}\\
$\calR$; $L_m$; $S_F$ & realignment \eqref{eq:realignment-kernel}; fixed corner kernel \eqref{eq:correction-cancellation}; tangent kernel \eqref{eq:tangent-kernel}\\
$\xi,\eta,\zeta,\omega$ & complex kernel arguments\\
$f_{pq}(\xi,\eta)$; $\mathbf e_{pq}$ & pure feature; mixed basis vector\\
$A,C,\Delta$; $\Pi$ & kernel denominators; product of the four arguments (Section~\ref{sec:tails})\\
$k$, $\epsilon=2^{-k}$, $m=\epsilon^{-2}$ & dyadic exponent, scale and corner depth\\
$a_\ell,b_\ell$; $c_\ell$; $\hat c_j$ & witness speeds; witness coefficients; their integer weights \eqref{eq:witness}\\
$\mathsf G_m$; $\Phi_m$ & evaluated Gram matrix; its pairing with $c$ \eqref{eq:nodes}\\
$\alpha,\beta$; $\mathcal H$ & continuation coordinates (Section~\ref{sec:continuation}); diagonal-ray integral \eqref{eq:energy}\\
\bottomrule
\end{tabular}
\end{center}

\emph{Kernel conventions.} A Gram matrix $G$ represents the polynomial
\(q^TGq\), so a symmetric off-diagonal entry \(G_{ij}\) contributes
\(2G_{ij}q_iq_j\); every matrix below is normalized this way. Pure
wedges are evaluated by the feature
\(f_{pq}(\xi,\eta)=\xi^p\eta^q-\xi^q\eta^p\), mixed wedges by
\(\xi^p\eta^q\). A matrix $G$ is identified with its four-slot
generating kernel $G(\xi,\eta,\zeta,\omega)=f(\xi,\eta)^TGf(\zeta,\omega)$,
formed with the appropriate features, and
\(P[\xi,\eta]=P(\xi,\eta,\bar \xi,\bar \eta)\) denotes the Hermitian
diagonal of a kernel $P$. Sums of pure and mixed matrices are sums of
their kernels. A subscript $m$ denotes a finite depth-$m$ kernel, and a
subscript $\infty$ its absolutely convergent infinite series.

\emph{Realignment.} For a pure kernel $P$, the mixed kernel $\calR P$ is
\begin{equation}
 (\mathcal RP)(\xi,\eta,\zeta,\omega)=P(\xi,\zeta,\eta,\omega).                    \label{eq:realignment-kernel}
\end{equation}
In entries, extend a symmetric pure matrix $\Xi$ antisymmetrically in
each index pair, $\Xi_{qp,rs}=\Xi_{pq,sr}=-\Xi_{pq,rs}$ and
$\Xi_{pp,rs}=0$. Then $\calR\Xi$ is the symmetric mixed matrix with
\begin{equation}\label{eq:realignment-entries}
 \calR(\Xi)_{pr,qs}=\Xi_{pq,rs},\qquad
 \calR(\Xi)_{ps,qr}=-\Xi_{pq,rs}\qquad(p<q,\ r<s),
\end{equation}
and every other entry is a symmetric counterpart of these or zero.
Indeed
$\Xi(\xi,\zeta,\eta,\omega)=\sum_{p<q,\,r<s}\Xi_{pq,rs}
(\xi^p\zeta^q-\xi^q\zeta^p)(\eta^r\omega^s-\eta^s\omega^r)$, and the
mixed monomials $\xi^p\eta^r\zeta^q\omega^s$ and
$\xi^p\eta^s\zeta^q\omega^r$ carry the entries $(pr,qs)$ and $(ps,qr)$.

A \emph{pure four-form} is a pure matrix $E_{pq,rs}=\Omega_{pqrs}$
($p<q$, $r<s$) with $\Omega$ totally antisymmetric in its four indices;
these are exactly the combinations of the pure Pl\"ucker matrices of
Lemma~\ref{lem:frame}. Its kernel
$E(\xi,\eta,\zeta,\omega)=\sum_{p,q,r,s}\Omega_{pqrs}\xi^p\eta^q\zeta^r\omega^s$
is alternating in its four arguments. Hence $\mathcal RE=-E$, and $E$
vanishes when two arguments coincide.

\subsection{The normal form}
By Lemma~\ref{lem:frame}, the Gram matrices of the corner form a coset of
the Pl\"ucker span. Two symmetries of the corner reduce this coset to
two corrections.

\begin{lemma}[Stabilization and the complete normal form]\label{lem:corner}
Let $N\ge2m$. The restriction of $F_N$ to the retained pairs
$w_p^\pm$, $0\le p<m$, all other symbol variables being zero, is a
polynomial $\calB_m$ that does not depend on $N$. In the ordered wedge
coordinates $(u,v,z^{\rm mix})$ it is represented by the Gram matrix with pure
block \(\left(\begin{smallmatrix}D&K\\K&D\end{smallmatrix}\right)\),
mixed block \(B\), and no pure--mixed entries, where
\begin{equation}
\begin{aligned}
 D_{pq,rs}&=2(p+1)(q+1)\delta_{pr}\delta_{qs},\quad p<q,\ r<s,\\
 K_{pq,rs}&=-2\mathbf1_{q-p=s-r}\min(p+1,r+1),\quad p<q,\ r<s,\\
 B_{pq,rs}&=2(p+1)(q+1)\delta_{pr}\delta_{qs}
 -2\mathbf1_{p+q=r+s}\min(p+1,q+1,r+1,s+1).
\end{aligned}                                                    \label{eq:corner-baselines}
\end{equation}
Every positive semidefinite Gram matrix representing $\calB_m$ becomes,
after averaging over the side sign change $w^-\mapsto-w^-$ and the side
interchange $w^-\leftrightarrow w^+$,
\begin{equation}\label{eq:corner-gram}
 Q_{\rm av}=\begin{pmatrix}
 D+E&K+\Theta&0\\
 K+\Theta&D+E&0\\
 0&0&B-\calR\Theta
 \end{pmatrix}\succeq0,
\end{equation}
where \(E\) is a pure four-form and \(\Theta\) an unrestricted real
symmetric matrix on the pure wedge space. The orthogonal change to
$(u+v)/\sqrt2,(u-v)/\sqrt2,z^{\rm mix}$ gives the positive semidefinite blocks
\begin{equation}
 P_+=D+E+K+\Theta,\qquad P_-=D+E-K-\Theta,\qquad M=B-\mathcal R\Theta.             \label{eq:corner-sectors}
\end{equation}
In particular \(\mathcal RE=-E\).
\end{lemma}
\begin{proof} \emph{Baseline.} Write $T_N(x)=\sum_\lambda x_\lambda S_\lambda$
with the shift matrices $S_\lambda$. Distinct shifts are Frobenius
orthogonal, and \(\|S_\lambda\|_F^2=N-|\lambda|\), which is \(p+1\) at
depth \(p\). The weighted Lagrange identity for the first two terms of
\eqref{eq:FN} and bilinearity of the commutator give
\[
 F_N=\sum_{\lambda<\lambda'}2(N-|\lambda|)(N-|\lambda'|)\,z_{\lambda\lambda'}^2
 -\Bigl\|\sum_{\lambda<\lambda'}z_{\lambda\lambda'}[S_\lambda,S_{\lambda'}]\Bigr\|_F^2 .
\]
On the retained pairs, same-side shifts commute, so only mixed wedges
enter the commutator term, and direct multiplication gives
\begin{equation}
 \langle[S_{\lambda_p^-},S_{\lambda_q^+}],[S_{\lambda_r^-},S_{\lambda_s^+}]\rangle_F
   =2\mathbf1_{p-q=r-s}\min(p+1,q+1,r+1,s+1).           \label{eq:corner-commutators}
\end{equation}
Each such commutator is supported on two strings of entries near
opposite corners of the matrix; for \(N\ge2m\) the strings are disjoint
and independent of \(N\). This proves the stabilization. The resulting
raw Gram matrix has the diagonal of $D$ and $B$, and its only
off-diagonal entries join mixed wedges $z^{\rm mix}_{pq},z^{\rm mix}_{rs}$ with $p-q=r-s$.
For such a pair with $p<r$ and $q<s$, the entry
$-2\min(p+1,q+1)$ contributes $-4\min(p+1,q+1)\,z^{\rm mix}_{pq}z^{\rm mix}_{rs}$. The
Pl\"ucker relation
\begin{equation}\label{eq:plucker-exchange}
 u_{pr}v_{qs}-z^{\rm mix}_{pq}z^{\rm mix}_{rs}+z^{\rm mix}_{ps}z^{\rm mix}_{rq}=0
\end{equation}
rewrites this contribution as
$-4\min(p+1,q+1)(u_{pr}v_{qs}+z^{\rm mix}_{ps}z^{\rm mix}_{rq})$: a pure pair with equal gaps
$r-p=s-q$ and a mixed pair with equal totals $p+s=r+q$. Summing over
all raw pairs gives exactly $K$ and $B$ in \eqref{eq:corner-baselines}.

\emph{Freedom.} By Lemma~\ref{lem:frame}, any other Gram matrix of
$\calB_m$ differs from this one by a combination of Pl\"ucker matrices on
four retained pairs. Four pairs from one side give a pure relation in
$u$ or in $v$; three pairs from one side give a relation coupling pure
and mixed wedges; two pairs from each side give
\eqref{eq:plucker-exchange} up to relabelling. Adding $2\theta$ times
\eqref{eq:plucker-exchange} adds $\theta$ to the cross entry
$(u_{pr},v_{qs})$, and $-\theta$ and $\theta$ to the mixed entries
$(z^{\rm mix}_{pq},z^{\rm mix}_{rs})$ and $(z^{\rm mix}_{ps},z^{\rm mix}_{rq})$. By
\eqref{eq:realignment-entries}, a cross correction $\Theta$ thus comes
with the mixed correction $-\calR\Theta$.

\emph{Averaging.} The side sign change preserves $\calB_m$, because the
baseline has no pure--mixed entries; the side interchange transposes
both Toeplitz matrices and preserves \eqref{eq:FN}. Conjugating a positive semidefinite
representing Gram matrix by either map gives another one. The sign
change fixes $u,v$ and negates $z^{\rm mix}$, so averaging over it removes all
pure--mixed couplings. The interchange maps $u_{pq}\mapsto v_{pq}$ and
$z^{\rm mix}_{pq}\mapsto-z^{\rm mix}_{qp}$, so averaging over it makes the two pure
diagonal blocks equal, with a common four-form $E$, and the cross block
symmetric. This gives \eqref{eq:corner-gram}; the orthogonal change of
coordinates gives \eqref{eq:corner-sectors}.
\end{proof}

Set \(Q_0=D+K\). An exact identity now removes both arbitrary corrections:
\begin{equation}
 \boxed{\ M+(P_++P_-)/2+\mathcal RP_+
                       =B+D+\mathcal RQ_0=:L_m.\ }     \label{eq:correction-cancellation}
\end{equation}
Indeed the correction terms are \(-\mathcal R\Theta+E+\mathcal RE+\mathcal R\Theta=0\).
This algebra uses no positivity. Positivity is used only for the first
three terms, which are kernels of positive semidefinite matrices; the
realigned term need not be positive. This is the identity
\eqref{eq:mechanism-cancellation}.
